% ===========================================================================
%  sviluppatore/06-dashboard/02-dati.tex
% ===========================================================================
\section{Come i dati arrivano alle pagine}\label{sez:data-layer}

La dashboard non legge i \file{parquet} a ogni richiesta. All'avvio li carica
tutti in memoria, in un oggetto che viene poi passato a tutto il resto:

\begin{codiceInterno}{python}{L'interfaccia pubblica}{dashboard/data\_layer.py}
    DataLayer()
        .startup()                  - carica streams + computa ensembles
        .get_ensemble(name)         - 'core' | 'full'
        .get_stream(name)           - stream singolo
        .get_stream_names()         - lista
        .get_spy_benchmark()        - SPY returns
\end{codiceInterno}

\noindent
È il secondo argomento che ogni modulo \file{ui\_ext\_*} riceve in
\simbolo{register(app, data\_layer, base\_path)}, e il motivo per cui una
pagina nuova non deve aprire file.

\subsection{Che cosa è pigro e che cosa no}

Gli stream si caricano all'avvio; due serie costose no.

\begin{codiceInterno}{python}{Due eccezioni dichiarate}{dashboard/data\_layer.py:86}
        # Stress di mercato (iso_stress) e risk-off macro (macro_regime): caricati o
        # CALCOLATI ON-DEMAND e cacheati (lazy), solo quando l'overlay corrispondente e'
        # attivo -> vedi _get_market_stress / _get_macro_riskoff. Nessun calcolo a startup.
\end{codiceInterno}

\noindent
Il criterio: una serie che serve solo a un overlay si calcola quando
quell'overlay è attivo. Caricarla all'avvio allungherebbe la partenza per
tutti, a beneficio di chi usa una funzione fra le sette.

\subsection{Gli ensemble preconfigurati non esistono più}

\begin{codiceInterno}{python}{La decisione, scritta nel codice}{dashboard/data\_layer.py:100}
        # Ensemble pre-configurati DISABILITATI: il sistema usa solo l'ensemble
        # custom (pagina /ensemble_custom), costruito a runtime dagli stream
        # selezionati. CORE/FULL rimossi su richiesta. Set vuoto => nessuna
        # scheda ENSEMBLE fissa in nav/home (restano stream singoli + custom).
        self.configured_ensemble_names: set[str] = set()
\end{codiceInterno}

\noindent
Il metodo \cod{get\_ensemble(name)} esiste ancora e il suo docstring parla
ancora di \cod{'core'} e \cod{'full'}: sono nomi che non si possono più
ottenere. Chi cerca di capire perché la pagina \percorsoHTTP{/ensemble/core} dà
\cod{404} lo trova qui.

\subsection{Come uno stream entra nell'elenco}

\file{dashboard/config.py} costruisce l'elenco degli stream in tre passi, e
solo il primo è un elenco scritto a mano:

\begin{enumerate}
  \item Gli stream dichiarati esplicitamente nella mappa interna.
  \item \textbf{Le strategie scoperte}, tramite
        \cod{discover\_strategies()}: entra chi ha un
        \cod{output\_path\_pattern}; chi non ce l'ha viene saltato con un
        \cod{continue} e un commento che dice «strategia non self-contained
        per dashboard».
  \item \textbf{Gli approcci scoperti}, con lo stesso meccanismo.
\end{enumerate}

\begin{perche}{Perché il passo 3 usa \cod{discover\_all()} e non \cod{discover\_approaches()}}
Il commento lo spiega: «anche gli overlay producono uno stream, e uno stream
va elencato indipendentemente da chi lo ha generato. Escluderli qui li
farebbe sparire dalle composizioni, che è l'unico posto in cui servono
davvero». È il complemento della scelta fatta nel registro
(\vediSez{nuovo-approccio}): gli overlay restano \emph{fuori} dall'elenco dei
backtest lanciabili e \emph{dentro} l'elenco degli stream componibili. Due
elenchi, due domande diverse.
\end{perche}

\begin{attenzione}{Il fallimento della scoperta è un avviso}
Entrambi i passi di scoperta sono dentro un \cod{try/except} che registra un
avviso e prosegue. Un errore di import in una strategia non impedisce alla
dashboard di partire: fa sparire quella strategia dall'elenco, e lascia una
riga \cod{Self-contained stream discovery failed} nel registro. È lo stesso
schema delle estensioni v2, con la stessa conseguenza: si nota guardando
l'interfaccia, non guardando un errore.
\end{attenzione}
